% --- 1. INTRODUZIONE ---
\section{Introduzione}

\subsection{Scopo del documento}
[-]

\subsection{Criteri di valutazione}
Per garantire una scelta oggettiva, il gruppo ha definito i seguenti criteri di valutazione. A ciascun criterio è stato assegnato un punteggio da 1 (pessimo) a 5 (ottimo):
\begin{itemize}
	\item \textbf{C1 - Interesse tecnico:} {[-]}
	\item \textbf{C2 - Complessità:} {[-]}
	\item \textbf{C3 - Chiarezza requisiti:} {[-]}
	\item \textbf{C4 - Qualità rapporto proponente:} {[-]}
	\item \textbf{C5 - Fattibilità tecnica:} {[-]}
	\item \textbf{C6 - Fattibilità economica:} {[-]}
	\item \textbf{C7 - Allineamento competenze:} {[-]}
\end{itemize}

\newpage
% --- 2. CONFRONTO CAPITOLATI ---
\section{Confronto Capitolati}
La seguente tabella riassume i punteggi (da 1 a 5) assegnati dal gruppo a ciascun capitolato analizzato, basandosi sulla presentazione iniziale, sullo studio dei documenti forniti, sugli scambi email con i proponenti e sull'analisi interna.

\vspace{0.5cm}
\begin{table}[ht]
	\centering
	\renewcommand{\arraystretch}{1.3}
	\begin{tabularx}{\textwidth}{X c c c c c c c | c}
		\toprule
		\textbf{Capitolato} & \textbf{C1} & \textbf{C2} & \textbf{C3} & \textbf{C4} & \textbf{C5} & \textbf{C6} & \textbf{C7} & \textbf{Totale} \\
		\midrule
		Capitolato x: {[Nome]} & [Num] & [Num] & [Num] & [Num] & [Num] & [Num] & [Num] & \textbf{[Tot]} \\
		Capitolato x: {[Nome]} & [Num] & [Num] & [Num] & [Num] & [Num] & [Num] & [Num] & \textbf{[Tot]} \\
		Capitolato x: {[Nome]} & [Num] & [Num] & [Num] & [Num] & [Num] & [Num] & [Num] & \textbf{[Tot]} \\
		\bottomrule
	\end{tabularx}
\end{table}
\vspace{0.2cm}
\textit{*Legenda Criteri: C1=Interesse, C2=Complessità, C3=Chiarezza req., C4=Rapporto prop., C5=Fatt. tecnica, C6=Fatt. economica, C7=Allineamento.}

\vspace{1cm}